% AA-11 | Algebra abstrakcyjna | Iteracja VI
\chapter{Pierścienie wielomianów i kryteria nierozkładalności}
\chaptermark{Wielomiany i nierozkładalność}
\label{aa:ch:wielomiany}

\section{Dlaczego rozróżniamy wielomian i funkcję wielomianową?}
Wielomian jest formalną skończoną sumą współczynników i potęg
symbolu $x$. Jego wartość w danym elemencie pierścienia otrzymujemy
przez podstawienie, lecz wielomian nie jest z definicji funkcją.
To rozróżnienie jest ważne już nad ciałami skończonymi.

\begin{definicja}[Pierścień wielomianów jednej zmiennej]
\label{aa:def:poly-ring}
Niech $R$ będzie przemiennym pierścieniem z jedynką.
Elementem $R[x]$ jest ciąg $(a_0,a_1,\ldots)$ elementów $R$,
w którym tylko skończenie wiele współczynników jest niezerowych.
Zapisujemy go jako $f=\sum_{j=0}^n a_jx^j$.
Dodawanie jest współczynnikowe, natomiast współczynnik przy $x^k$
iloczynu $fg$ jest równy $\sum_{i+j=k}a_i b_j$.
Wielomian zerowy ma wszystkie współczynniki równe zeru.
Jeśli $f\ne0$, jego \emph{stopniem} jest największy indeks
niezerowego współczynnika, a \emph{współczynnikiem wiodącym}
--- współczynnik o tym indeksie. Wielomian jest \emph{unormowany},
gdy jego współczynnik wiodący jest równy $1_R$.
\end{definicja}

\begin{lemat}[Aksjomaty pierścienia wielomianów]
\label{aa:lem:poly-ring-axioms}
Dla każdego przemiennego pierścienia $R$ z jedynką
wzory podane w Definicji~\ref{aa:def:poly-ring}
wyznaczają przemienny pierścień $R[x]$ z jedynką.
Odwzorowanie $a\mapsto(a,0,0,\ldots)$ jest zanurzeniem
pierścieni $R\to R[x]$ (także dla pierścienia zerowego).
\end{lemat}
\begin{proof}
Suma dwóch ciągów o skończonym nośniku ma skończony
nośnik. Jeżeli któryś z czynników jest wielomianem
zerowym, ich iloczyn jest zerowy. W przeciwnym razie
współczynniki iloczynu przy $x^k$ dla $k$ większego
od sumy największych niezerowych indeksów obu czynników
są zerowe. Zatem iloczyn także ma skończony nośnik.
Aksjomaty dodawania wynikają współczynnikowo z aksjomatów
$(R,+)$; elementem zerowym jest ciąg zerowy.
Przy mnożeniu przemienność wynika z wymiany indeksów
$i,j$ w sumie $\sum_{i+j=k}a_ib_j$. Przy łączności
współczynnik przy $x^k$ iloczynu $(fg)h$ ma postać
$\sum_{i+j+\ell=k} a_i b_j c_\ell$, tak samo
jak współczynnik przy $x^k$ iloczynu $f(gh)$.
Wszystkie sumy są skończone, zatem wolno zmienić kolejność
sumowania. Ciąg $(1_R,0,\ldots)$ jest jedynką,
a równość $f(g+h)=fg+fh$ można sprawdzić przy każdym
współczynniku, stosując rozdzielność w $R$.
Stałe wielomiany zachowują oba działania i jedynkę.
Jeśli dwa stałe wielomiany są równe, ich współczynniki
przy $x^0$ są równe, więc odwzorowanie jest różnowartościowe
również dla pierścienia zerowego.
\end{proof}

\begin{lemat}[Stopień iloczynu nad dziedziną]
\label{aa:lem:poly-degree}
Niech $R$ będzie dziedziną całkowitości i niech
$f,g\in R[x]$ będą niezerowe. Wówczas
$\deg(fg)=\deg f+\deg g$. Jednostkami pierścienia $R[x]$
są dokładnie wielomiany stałe, których jedyny współczynnik
jest jednostką w $R$.
\end{lemat}
\begin{proof}
Niech $n=\deg f$, $m=\deg g$ oraz $a_n,b_m$ będą
współczynnikami wiodącymi. Przy potędze $x^{n+m}$ w iloczynie
może wystąpić tylko składnik $a_nb_m$. Jest on niezerowy,
ponieważ $R$ nie ma dzielników zera. Wyższe współczynniki są
zerowe, gdyż nie istnieje para indeksów $i\le n$, $j\le m$
z sumą większą niż $n+m$. To dowodzi wzoru na stopień.
Jeśli $fg=1_R$, obydwa czynniki są niezerowe, a suma
ich stopni jest zerem. Oba stopnie wynoszą zatem zero,
więc $f,g$ są stałe i ich współczynniki są wzajemnymi
odwrotnościami. Odwrotnie, każda stała jednostka $R$
jest jednostką $R[x]$.
\end{proof}

\begin{definicja}[Wartość wielomianu]
\label{aa:def:poly-evaluation}
Niech $R$ będzie przemiennym pierścieniem z jedynką,
$a\in R$ oraz $f=\sum_{j=0}^n c_jx^j\in R[x]$.
Definiujemy
$f(a)=\sum_{j=0}^n c_ja^j$. Odwzorowanie
$\operatorname{ev}_a:R[x]\to R$, $f\mapsto f(a)$,
nazywamy \emph{homomorfizmem podstawienia}.
Jest on rzeczywiście homomorfizmem: wzór na współczynniki
iloczynu i rozdzielność w $R$ dają
$(fg)(a)=f(a)g(a)$, a zgodność z dodawaniem i jedynką
wynika bezpośrednio z definicji.
\end{definicja}
\begin{lemat}[Własność uniwersalna pierścienia wielomianów]
\label{aa:lem:poly-universal}
Niech $R,S$ będą przemiennymi pierścieniami z jedynką,
$\varphi:R\to S$ homomorfizmem pierścieni i $s\in S$.
Istnieje dokładnie jeden homomorfizm pierścieni
$\widetilde\varphi:R[x]\to S$, którego ograniczenie
do $R$ jest równe $\varphi$ oraz dla którego
$\widetilde\varphi(x)=s$. Ma on postać
$\widetilde\varphi(\sum_j a_jx^j)=\sum_j\varphi(a_j)s^j$.
\end{lemat}
\begin{proof}
Homomorfizm musi spełniać podany wzór, ponieważ zachowuje
sumy, iloczyny i jedynkę. Wzór jest dobrze określony:
każdy wielomian ma dokładnie jeden skończony zestaw współczynników.
Zachowanie dodawania wynika ze współczynnikowego dodawania.
Dla mnożenia mamy, z przemienności w $S$,
\[
 \sum_k\varphi\Bigl(\sum_{i+j=k}a_i b_j\Bigr)s^k
 =\sum_{i,j}\varphi(a_i)s^i\varphi(b_j)s^j,
\]
co jest iloczynem obrazów dwóch wielomianów.
Obraz stałej $1_R$ jest $1_S$. Otrzymaliśmy homomorfizm;
jego jednoznaczność została wykazana na początku.
\end{proof}

\begin{przyklad}[Dwa różne wielomiany wyznaczające tę samą funkcję]
\label{aa:prz:poly-finite-field}
W ciele $\mathbb F_p=\mathbb Z/p\mathbb Z$,
gdzie $p$ jest liczbą pierwszą, wielomian $x^p-x$
nie jest wielomianem zerowym: ma dwa niezerowe współczynniki.
Jednak dla każdej klasy $a\in\mathbb F_p$ zachodzi
$a^p=a$, na mocy małego twierdzenia Fermata dla $a\ne0$
i bezpośrednio dla $a=0$. Funkcja zadana przez $x^p-x$
jest więc identycznie zerowa. Tym samym różne wielomiany
nie muszą wyznaczać różnych funkcji na ciele skończonym.
\end{przyklad}



\section{Dzielenie z resztą, pierwiastki i czynniki liniowe}
\begin{twierdzenie}[Dzielenie wielomianów przez wielomian o współczynniku wiodącym odwracalnym]
\label{aa:tw:poly-division}
Niech $R$ będzie przemiennym pierścieniem z jedynką,
$f,g\in R[x]$, $g\ne0$, a współczynnik wiodący $g$
niech będzie jednostką $R$. Istnieją dokładnie jedne
$q,r\in R[x]$ takie, że $f=qg+r$ oraz $r=0$ albo
$\deg r<\deg g$. W szczególności twierdzenie zachodzi
dla dowolnego $g\ne0$ nad ciałem $K$.
\end{twierdzenie}
\begin{proof}
Niech $d=\deg g$ i niech $u$ będzie odwracalnym
współczynnikiem wiodącym $g$.
Zaczynamy od $f_0=f$ i $q_0=0$. Jeśli $f_0=0$
lub $\deg f_0<d$, kończymy. W przeciwnym razie,
przy $\deg f_0=n$ i współczynniku wiodącym $a$,
odejmujemy wielomian $au^{-1}x^{n-d}g$ i dopisujemy
$au^{-1}x^{n-d}$ do ilorazu. Współczynnik przy
$x^n$ znika, więc stopień nowej reszty jest mniejszy
od $n$. Procedura kończy się po skończonej liczbie kroków,
bo stopnie niezerowych reszt są nieujemnymi liczbami całkowitymi
i ściśle maleją. Sumując wykonane odejmowania,
otrzymujemy $f=qg+r$ z wymaganym ograniczeniem.

Dla jednoznaczności załóżmy, że
$f=qg+r=q'g+r'$ z $\deg r,\deg r'<d$ dla niezerowych
reszt. Wtedy $(q-q')g=r'-r$. Jeśli $q\ne q'$,
wyraz wiodący $(q-q')g$ jest niezerowy: współczynnik
wiodący $g$ jest jednostką, więc nie może wyzerować
niezerowego współczynnika wiodącego $q-q'$.
Stąd $\deg((q-q')g)=\deg(q-q')+d\ge d$,
co jest sprzeczne z tym, że $r'-r$ jest wielomianem
zerowym albo ma stopień mniejszy niż $d$. Zatem
$q=q'$ i $r=r'$. 
\end{proof}

\begin{wniosek}[Twierdzenie o czynniku liniowym]
\label{aa:wn:poly-factor}
Jeśli $R$ jest przemiennym pierścieniem z jedynką,
$f\in R[x]$ i $a\in R$, to $f(a)=0$ wtedy i tylko wtedy,
gdy $x-a$ dzieli $f$ w $R[x]$.
\end{wniosek}
\begin{proof}
Jeśli $R$ jest pierścieniem zerowym, to $a=0_R$,
$f=0_{R[x]}$ i $x-a=0_{R[x]}$; zachodzą zarówno
$f(a)=0_R$, jak i $(x-a)\mid f$, gdyż $0=0\cdot0$. Załóżmy zatem,
że $R$ jest niezerowy. Wtedy $x-a$ jest niezerowym
wielomianem o współczynniku wiodącym $1_R$,
więc poprzednie twierdzenie daje $f=(x-a)q+r$,
gdzie $r$ jest stałą. Po podstawieniu $x=a$ dostajemy
$f(a)=r$. Zatem $f(a)=0$ jest równoważne $r=0$,
czyli podzielności przez $x-a$.
\end{proof}

\begin{twierdzenie}[Ograniczenie liczby pierwiastków]
\label{aa:tw:poly-roots}
Jeżeli $R$ jest dziedziną całkowitości, a
$f\in R[x]$ niezerowym wielomianem stopnia $n$, to
$f$ ma co najwyżej $n$ różnych pierwiastków w $R$.
\end{twierdzenie}
\begin{proof}
Indukujemy po $n$. Dla $n=0$ niezerowa stała
nie ma pierwiastków. Niech $n\ge1$. Jeśli $f$ nie ma
pierwiastków, twierdzenie zachodzi. W przeciwnym razie
wybierzmy pierwiastek $a$. Z wniosku poprzedniego wynika
$f=(x-a)g$, a porównanie stopni daje $\deg g=n-1$.
Jeśli $b$ jest innym pierwiastkiem $f$, to
$0=f(b)=(b-a)g(b)$. Ponieważ $b-a\ne0$ i $R$ nie
ma dzielników zera, otrzymujemy $g(b)=0$.
Z założenia indukcyjnego wielomian $g$ ma najwyżej
$n-1$ różnych pierwiastków, a dodanie $a$ daje
co najwyżej $n$ pierwiastków $f$.
\end{proof}

\begin{twierdzenie}[Funkcje wielomianowe na ciele skończonym]
\label{aa:tw:finite-field-polynomial-functions}
Dla liczby pierwszej $p$ każda funkcja
$F:\mathbb F_p\to\mathbb F_p$ jest zadana przez
dokładnie jeden wielomian zerowy lub stopnia mniejszego niż $p$.
Jądro homomorfizmu
\[
 \Theta:\mathbb F_p[x]\longrightarrow
 \{F:\mathbb F_p\to\mathbb F_p\},\qquad
 \Theta(f)(a)=f(a),
\]
gdzie po prawej działania wykonuje się punktowo,
jest ideałem $(x^p-x)$.
\end{twierdzenie}
\begin{proof}
Dla $a\in\mathbb F_p$ rozważmy wielomian Lagrange'a
$L_a(x)=\prod_{b\in\mathbb F_p,\,b\ne a}
 (x-b)/(a-b)$. Wszystkie mianowniki są niezerowe,
więc odwracalne. Mamy $L_a(a)=1$ oraz $L_a(b)=0$
dla $b\ne a$. Wielomian
$P_F(x)=\sum_{a\in\mathbb F_p}F(a)L_a(x)$
ma stopień co najwyżej $p-1$ i spełnia
$P_F(b)=F(b)$ dla każdego $b\in\mathbb F_p$.

Jeżeli dwa wielomiany stopnia mniejszego niż $p$
wyznaczają tę samą funkcję, ich różnica ma co najmniej
$p$ różnych pierwiastków. Na mocy
Twierdzenia~\ref{aa:tw:poly-roots} ta różnica musi być
wielomianem zerowym.

Wielomian $x^p-x$ zeruje się we wszystkich punktach,
więc jego wszystkie wielokrotności należą do jądra.
Dla dowolnego $f$ dzielenie z resztą przez unormowany
$x^p-x$ daje $f=q(x^p-x)+r$ z $r=0$ lub
$\deg r<p$. Jeśli $f$ jest w jądrze, to $r$
zeruje się we wszystkich $p$ punktach, a więc $r=0$.
Odwrotne zawieranie jest oczywiste, co kończy dowód.
\end{proof}

\begin{przyklad}[Dlaczego nie można opuścić założenia o dziedzinie]
\label{aa:prz:poly-roots-z8}
W $\mathbb Z/8\mathbb Z$ wielomian $x^2-1$
ma cztery różne pierwiastki: $[1],[3],[5],[7]$.
Nie przeczy to twierdzeniu, ponieważ pierścień
$\mathbb Z/8\mathbb Z$ nie jest dziedziną całkowitości.
Różnica wynika z zerowania iloczynu $(b-a)g(b)$
przy obu niezerowych czynnikach.
\end{przyklad}

\begin{wniosek}[Wielomian nad nieskończoną dziedziną jest określony przez wartości]
\label{aa:wn:poly-evaluation-infinite-domain}
Niech $R$ będzie nieskończoną dziedziną całkowitości.
Jeśli $f,g\in R[x]$ spełniają $f(a)=g(a)$ dla każdego
$a\in R$, to $f=g$. Równoważnie, odwzorowanie przypisujące
wielomianowi jego funkcję wartości na $R$ jest różnowartościowe.
\end{wniosek}
\begin{proof}
Połóżmy $h=f-g$. Z założenia $h(a)=0$ dla każdego
$a\in R$. Gdyby $h\ne0$, to wielomian $h$ miałby
skończony stopień $n$ i, na mocy
Twierdzenia~\ref{aa:tw:poly-roots}, najwyżej $n$
różnych pierwiastków w $R$. Tymczasem każdy element
nieskończonego zbioru $R$ byłby jego pierwiastkiem.
Otrzymana sprzeczność daje $h=0$, czyli $f=g$.
\end{proof}
\begin{uwaga}
Założenia o nieskończoności i braku dzielników zera
pełnią różne funkcje. Nad ciałem skończonym odwzorowanie
wartościowania może nie być różnowartościowe
(Przykład~\ref{aa:prz:poly-finite-field}); w pierścieniu
z dzielnikami zera zawodzi oszacowanie liczby pierwiastków
(Przykład~\ref{aa:prz:poly-roots-z8}).
\end{uwaga}

\section{Wielomiany pierwotne i lemat Gaussa}
\begin{definicja}[Ciało ułamków]
\label{aa:def:fracfield}
Niech $R$ będzie dziedziną całkowitości. Jej
\emph{ciałem ułamków} jest ciało $K=\operatorname{Frac}(R)$
przedstawione przez klasy par $(a,b)$, gdzie $a\in R$
i $b\in R\setminus\{0\}$, przy relacji
$(a,b)\sim(c,d)$ wtedy i tylko wtedy, gdy $ad=bc$.
Klasę pary $(a,b)$ oznaczamy $a/b$.
Dodawanie i mnożenie ułamków wyznaczają wzory
$a/b+c/d=(ad+bc)/(bd)$ oraz $(a/b)(c/d)=ac/(bd)$.
Odwzorowanie $R\to K$, $a\mapsto a/1$, jest zanurzeniem.
\end{definicja}

\begin{lemat}[Uzasadnienie konstrukcji ciała ułamków]
\label{aa:lem:fracfield-well-defined}
Niech $R$ będzie dziedziną całkowitości. Relacja
$(a,b)\sim(c,d)$, określona przez $ad=bc$ na
$R\times(R\setminus\{0\})$, jest relacją równoważności.
Wzory na dodawanie i mnożenie ułamków nie zależą od
wyboru przedstawicieli, a zbiór klas z tymi działaniami
jest ciałem zawierającym izomorficzną kopię $R$.
\end{lemat}
\begin{proof}
Zwrotność i symetria są bezpośrednie. Jeśli
$ad=bc$ oraz $cf=de$, to $adf=bcf=bde$.
Ponieważ $d\ne0$ i $R$ nie ma dzielników zera,
możemy skrócić $d$, otrzymując $af=be$;
to dowodzi przechodniości.

Sprawdźmy niezależność dodawania od przedstawicieli.
Załóżmy $ab'=a'b$ i $cd'=c'd$ z niezerowymi
$b,b',d,d'$. Wówczas
\[
 (ad+bc)b'd'=ab'dd'+bb'cd'
             =a'bd'd+bb'c'd
             =(a'd'+b'c')bd.
\]
Stąd $(ad+bc)/(bd)=(a'd'+b'c')/(b'd')$.
Dla mnożenia odpowiednia równość ma postać
$acb'd'=(ab')(cd')=(a'b)(c'd)=a'c'bd$.
Mianowniki $bd$ są niezerowe, gdyż $R$ jest dziedziną.

Łączność, przemienność i rozdzielność sprowadzają się
po rozwinięciu liczników do odpowiednich własności $R$.
Dla przykładu
$(a/b+c/d)+e/f=(adf+bcf+ebd)/(bdf)$, niezależnie
od nawiasowania. Zerem jest $0/1$, jedynką $1/1$,
a przeciwnym do $a/b$ element $(-a)/b$.
Z równości $a/b=0/1$ wynika $a=0$.
Jeśli $a\ne0$, to $b/a$ jest poprawnym ułamkiem
i $(a/b)(b/a)=1/1$. Zatem otrzymujemy ciało.
Ponadto $a\mapsto a/1$ zachowuje działania i jest
różnowartościowe, bo $a/1=c/1$ wymusza $a=c$.
Dalsze własności tej konstrukcji omawiamy w
dalszej części skryptu.
\end{proof}

\begin{definicja}[Zawartość wielomianu]
\label{aa:def:poly-content}
Niech $R$ będzie dziedziną z jednoznacznym rozkładem
na czynniki i niech $f=\sum_{j=0}^n a_jx^j\in R[x]$
będzie niezerowy. \emph{Zawartością} $f$ nazywamy
największy wspólny dzielnik jego niezerowych współczynników,
wyznaczony z dokładnością do stowarzyszenia; oznaczamy go
$c(f)$. Wielomian $f$ jest \emph{pierwotny}, gdy $c(f)$
jest jednostką, czyli żaden element pierwszy nie dzieli
wszystkich współczynników $f$. W szczególności
$f=c(f)f_0$ dla pewnego pierwotnego $f_0\in R[x]$.
\end{definicja}

\begin{lemat}[Iloczyn wielomianów pierwotnych]
\label{aa:lem:gauss-primitive}
Niech $R$ będzie dziedziną z jednoznacznym rozkładem
na czynniki. Jeżeli $f,g\in R[x]$ są pierwotne,
to $fg$ jest pierwotny.
\end{lemat}
\begin{proof}
Przypuśćmy, że $fg$ nie jest pierwotny. Wtedy
istnieje element pierwszy $p\in R$ dzielący każdy
współczynnik iloczynu. Ponieważ $p$ jest pierwszy,
ideał $(p)$ jest pierwszy: jeśli $ab\in(p)$, to
$p\mid ab$, zatem $p\mid a$ lub $p\mid b$.
Na mocy kryterium z rozdziału~\ref{aa:ch:pierscienie}
iloraz $R/(p)$ jest dziedziną całkowitości.
Zredukujmy współczynniki $f,g$ modulo $(p)$.
Otrzymamy wielomiany $\bar f,\bar g\in(R/(p))[x]$.
Pierwotność oznacza, że przynajmniej jeden współczynnik
każdego z nich nie jest podzielny przez $p$,
więc $\bar f\ne0$ i $\bar g\ne0$.
Ale wszystkie współczynniki $fg$ są podzielne przez $p$,
więc $\bar f\bar g=0$. To niemożliwe w pierścieniu
wielomianów nad dziedziną całkowitości na mocy
Lematu~\ref{aa:lem:poly-degree}. Otrzymujemy sprzeczność.
\end{proof}

\begin{twierdzenie}[Lemat Gaussa o zawartości]
\label{aa:tw:gauss-content}
Niech $R$ będzie dziedziną z jednoznacznym rozkładem
na czynniki i niech $f,g\in R[x]$ będą niezerowe.
Wówczas $c(fg)\sim c(f)c(g)$, gdzie $\sim$ oznacza
stowarzyszenie w $R$.
\end{twierdzenie}
\begin{proof}
Wybierzmy zawartości $c(f),c(g)$ i zapiszmy
$f=c(f)f_0$, $g=c(g)g_0$ z pierwotnymi $f_0,g_0$.
Wtedy $fg=c(f)c(g)(f_0g_0)$.
Poprzedni lemat mówi, że $f_0g_0$ jest pierwotny.
Zatem największy wspólny dzielnik współczynników iloczynu
jest stowarzyszony z $c(f)c(g)$, co daje tezę.
\end{proof}

\begin{lemat}[Usuwanie mianowników z rozkładu]
\label{aa:lem:gauss-denominators}
Niech $R$ będzie dziedziną z jednoznacznym rozkładem
na czynniki, $K=\operatorname{Frac}(R)$, a
$f\in R[x]$ wielomianem pierwotnym. Jeżeli
$f=gh$ dla wielomianów $g,h\in K[x]$ o dodatnich stopniach,
to istnieją wielomiany $g_0,h_0\in R[x]$ o dodatnich
stopniach takie, że $f=g_0h_0$.
\end{lemat}
\begin{proof}
Każdy niezerowy wielomian $g\in K[x]$ można zapisać
jako $g=\alpha G$, gdzie $\alpha\in K^\times$,
a $G\in R[x]$ jest pierwotny: najpierw mnożymy przez
wspólny niezerowy mianownik współczynników, a potem
dzielimy uzyskany wielomian przez jego zawartość.
Podobnie $h=\beta H$ z pierwotnym $H\in R[x]$.
Zatem $f=\lambda GH$, gdzie $\lambda=\alpha\beta\in K^\times$,
a $GH$ jest pierwotny na mocy Lematu~\ref{aa:lem:gauss-primitive}.

Zapiszmy $\lambda=a/b$ z $a,b\in R$, $b\ne0$,
przy czym $a$ i $b$ nie mają wspólnego czynnika nierozkładalnego.
Z równości $bf=aGH$ wynika, że $b$ dzieli każdy
współczynnik wielomianu $aGH$. Jeśli $b$ nie byłby
jednostką, miałby czynnik pierwszy $p$. Wówczas $p$
dzieliłby każdy współczynnik $aGH$; ponieważ $p\nmid a$,
pierwszość $p$ wymuszałaby, że $p$ dzieli wszystkie
współczynniki $GH$, wbrew pierwotności $GH$.
Stąd $b$ jest jednostką, czyli $\lambda\in R$.
Jeśli $\lambda$ nie byłaby jednostką, dowolny jej
czynnik pierwszy dzieliłby wszystkie współczynniki $f$,
co przeczyłoby pierwotności $f$. Ostatecznie
$\lambda\in R^\times$ i $f=(\lambda G)H$ jest rozkładem
w $R[x]$ na dwa wielomiany dodatniego stopnia.
\end{proof}

\begin{lemat}[Pierwotne wielomiany stowarzyszone nad ciałem ułamków]
\label{aa:lem:primitive-associates}
Niech $R$ będzie dziedziną z jednoznacznym rozkładem
na czynniki, a $K=\operatorname{Frac}(R)$.
Jeżeli dwa niezerowe pierwotne wielomiany
$f,g\in R[x]$ spełniają $f=\lambda g$ dla
pewnego $\lambda\in K^\times$, to $\lambda\in R^\times$.
\end{lemat}
\begin{proof}
Wybierzmy $\lambda=a/b$, gdzie $a,b\in R$ są
niezerowe i nie mają wspólnego czynnika nierozkładalnego.
Równość $bf=ag$ zachodzi współczynnikowo.
Gdyby $b$ nie było jednostką, dowolny czynnik pierwszy
$p\mid b$ dzieliłby każdy współczynnik $ag$.
Ponieważ $p\nmid a$, każdy współczynnik $g$ byłby
podzielny przez $p$, wbrew pierwotności $g$.
Zatem $b\in R^\times$ i $\lambda\in R$.
Gdyby $a$ nie było jednostką, pewien czynnik pierwszy
$p\mid a$ dzieliłby wszystkie współczynniki $bf$,
a ponieważ $p\nmid b$, również wszystkie współczynniki
$f$, znów sprzecznie z pierwotnością.
Ostatecznie $a,b\in R^\times$, a więc $\lambda\in R^\times$.
\end{proof}

\begin{twierdzenie}[Kryterium nierozkładalności Gaussa]
\label{aa:tw:gauss-irreducibility}
Niech $R$ będzie dziedziną z jednoznacznym rozkładem
na czynniki, $K=\operatorname{Frac}(R)$, a
$f\in R[x]$ pierwotnym wielomianem dodatniego stopnia.
Wówczas $f$ jest nierozkładalny w $R[x]$ wtedy i tylko
wtedy, gdy jest nierozkładalny w $K[x]$.
\end{twierdzenie}
\begin{proof}
Jeżeli $f$ rozkłada się na dwa czynniki dodatniego stopnia
w $K[x]$, to poprzedni lemat daje taki sam rodzaj rozkładu
w $R[x]$, więc $f$ nie jest nierozkładalny w $R[x]$.

Odwrotnie, przypuśćmy, że $f=uv$ w $R[x]$, gdzie
$u,v$ są niejednostkami $R[x]$. Gdyby jeden z nich
był stały, byłby nieodwracalnym elementem $R$, który
dzieliłby wszystkie współczynniki $f$. To przeczy
pierwotności. Zatem $u$ i $v$ mają dodatnie stopnie,
więc stanowią właściwy rozkład $f$ także w $K[x]$.
W ten sposób dowiedliśmy obu implikacji.
\end{proof}

\begin{twierdzenie}[Twierdzenie Gaussa o pierścieniu wielomianów]
\label{aa:tw:ufd-polynomial}
Jeżeli $R$ jest dziedziną z jednoznacznym rozkładem
na czynniki, to pierścień $R[x]$ także jest dziedziną
z jednoznacznym rozkładem na czynniki. W konsekwencji,
jeżeli $K$ jest ciałem, to $K[x_1,\ldots,x_n]$
ma jednoznaczny rozkład na czynniki dla każdego $n\ge1$.
\end{twierdzenie}
\begin{proof}
Pierścień $R[x]$ jest dziedziną całkowitości na mocy
Lematu~\ref{aa:lem:poly-degree}. Ustalmy
$K=\operatorname{Frac}(R)$. Niezerowy wielomian
$f\in R[x]$ zapisujemy jako $f=c(f)f_0$,
gdzie $f_0$ jest pierwotny. Element $c(f)$ rozkłada się
w $R$ na czynniki nierozkładalne, jeśli nie jest jednostką.
Z kolei $K[x]$ jest dziedziną euklidesową
(Rozdział~\ref{aa:ch:podzielnosc}), a więc ma
jednoznaczny rozkład na czynniki. Rozłóżmy $f_0$ w $K[x]$
na wielomiany nierozkładalne, przy czym dla stałego
pierwotnego $f_0$ przyjmujemy pusty iloczyn i $r=0$.
Wybierając dla każdego czynnika pierwotnego reprezentanta
w $R[x]$, otrzymujemy $f_0=\lambda g_1\cdots g_r$
z pierwotnymi $g_i\in R[x]$ i $\lambda\in K^\times$.
Dla $r=0$ mamy $f_0=\lambda\in R^\times$.
Lemat~\ref{aa:lem:primitive-associates}, zastosowany do
pierwotnych $f_0$ oraz $g_1\cdots g_r$, pokazuje,
że $\lambda\in R^\times$.
Z kryterium nierozkładalności Gaussa każdy $g_i$ jest
nierozkładalny w $R[x]$. W ten sposób otrzymaliśmy
istnienie rozkładu na czynniki nierozkładalne.

Sprawdźmy jednoznaczność. W każdym rozkładzie na czynniki
nierozkładalne w $R[x]$ czynniki stałe odpowiadają
czynnikom nierozkładalnym zawartości $c(f)$; pozostałe
czynniki mają dodatnie stopnie i muszą być pierwotne,
gdyż niezerowa niejednostka dzieląca wszystkie ich
współczynniki dawałaby właściwy rozkład.
Takie czynniki są nierozkładalne również w $K[x]$.
Jednoznaczność rozkładu zawartości w $R$ i rozkładu
w $K[x]$ zapewnia sparowanie czynników.
Jeśli dwa pierwotne wielomiany różnią się mnożnikiem
$\lambda\in K^\times$, to na mocy
Lematu~\ref{aa:lem:primitive-associates} mamy
$\lambda\in R^\times$. Czynniki odpowiadają sobie
zatem z dokładnością do jednostek $R[x]$.

Wreszcie ciało $K$ samo jest dziedziną z jednoznacznym
rozkładem, a powtarzając wykazany wynik kolejno dla
$K[x_1]$, $K[x_1][x_2]$, $\ldots$ otrzymujemy tezę
o dowolnej skończonej liczbie zmiennych.
\end{proof}

\begin{przyklad}[Jednoznaczny rozkład bez główności wszystkich ideałów]
\label{aa:prz:ufd-not-pid}
Dla dowolnego ciała $K$ pierścień $K[x,y]$ ma
jednoznaczny rozkład na czynniki na mocy poprzedniego
twierdzenia. Jednak ideał $(x,y)$ nie jest główny:
wykazaliśmy to już szczegółowo w
Przykładzie~\ref{aa:prz:two-polynomials-not-pid}.
Oznacza to, że implikacja
„dziedzina z jednoznacznym rozkładem na czynniki
$\Rightarrow$ dziedzina ideałów głównych” jest fałszywa.
\end{przyklad}

\section{Kryterium Eisensteina i badanie redukcji modulo liczbę pierwszą}
\begin{twierdzenie}[Kryterium Eisensteina]
\label{aa:tw:eisenstein}
Niech $R$ będzie dziedziną z jednoznacznym rozkładem
na czynniki, $K=\operatorname{Frac}(R)$, a
$f=a_nx^n+\cdots+a_1x+a_0\in R[x]$ wielomianem
stopnia $n\ge1$. Załóżmy, że istnieje element pierwszy
$p\in R$ taki, że $p\nmid a_n$, $p\mid a_j$ dla
$0\le j<n$ oraz $p^2\nmid a_0$.
Wtedy $f$ jest nierozkładalny w $K[x]$.
\end{twierdzenie}
\begin{proof}
Niech $c=c(f)$ będzie wybraną zawartością $f$
i niech $f_0=f/c\in R[x]$. Z warunku
$p\nmid a_n$ wynika $p\nmid c$, a $f_0$ jest pierwotny.
Ponieważ $a_j=c\,b_j$, gdzie $b_j$ jest współczynnikiem
$f_0$, oraz $p\nmid c$, pierwszość $p$ pokazuje,
że $f_0$ spełnia wszystkie trzy warunki kryterium:
$p$ pozostaje dzielnikiem niższych współczynników,
nie dzieli współczynnika wiodącego, a $p^2$ nie dzieli
wyrazu wolnego, ponieważ $p\nmid c$.

Gdyby $f$ był rozkładalny w $K[x]$, byłby rozkładalny
także $f_0$. Z kryterium nierozkładalności Gaussa
otrzymalibyśmy $f_0=gh$ w $R[x]$, gdzie obydwa czynniki
mają dodatnie stopnie. Zredukujmy współczynniki modulo
ideał $(p)$. Ponieważ $R/(p)$ jest dziedziną,
redukcja daje równość
$\overline g\,\overline h=\overline{a_n/c}\,x^n$.
Współczynnik wiodący po prawej jest niezerowy,
a oba czynniki po lewej mają dodatnie stopnie:
ich współczynniki wiodące nie mogą być podzielne przez
$p$, gdyż ich iloczyn jest współczynnikiem wiodącym $f_0$.

W dziedzinie całkowitości każdy dzielnik monomianu
$c x^n$ z $c\ne0$ jest sam monomianem z niezerowym
współczynnikiem. Można to także zobaczyć bezpośrednio:
jeśli $\bar g$ miałby niezerowy wyraz niższego stopnia,
to jego najniższy niezerowy współczynnik, pomnożony
przez najniższy niezerowy współczynnik $\bar h$,
dałby niezerowy wyraz iloczynu o wykładniku mniejszym
niż $n$. Zatem $\bar g$ i $\bar h$ są monomianami
dodatnich stopni. Ich wyrazy wolne są więc zerowe,
co oznacza $p\mid g(0)$ i $p\mid h(0)$. Wtedy
$p^2\mid g(0)h(0)=f_0(0)$, sprzecznie z założeniem.
Stąd $f$ jest nierozkładalny w $K[x]$.
\end{proof}

\begin{przyklad}[Wielomiany $x^n-p$]
\label{aa:prz:xn-p}
Dla liczby pierwszej $p$ i $n\ge1$ wielomian
$x^n-p\in\mathbb Z[x]$ spełnia kryterium Eisensteina
z elementem pierwszym $p$. Współczynnik wiodący
wynosi $1$, wszystkie pozostałe współczynniki
są podzielne przez $p$, a $p^2$ nie dzieli wyrazu
wolnego $-p$. Zatem $x^n-p$ jest nierozkładalny
nad $\mathbb Q$. W szczególności $x^3-2$ jest
wielomianem minimalnym liczby $\sqrt[3]{2}$ nad $\mathbb Q$.
\end{przyklad}

\begin{przyklad}[Zmiana zmiennej ujawniająca kryterium]
\label{aa:prz:x4plus1}
Wielomian $f(x)=x^4+1$ nie spełnia wprost kryterium
Eisensteina z liczbą pierwszą $2$, ponieważ jego
wyraz wolny nie jest podzielny przez $2$. Rozważmy jednak
$f(x+1)=x^4+4x^3+6x^2+4x+2$.
Wszystkie współczynniki oprócz wiodącego są parzyste,
a $4$ nie dzieli wyrazu wolnego $2$.
Kryterium Eisensteina daje nierozkładalność $f(x+1)$
w $\mathbb Q[x]$. Odwzorowanie
$\mathbb Q[x]\to\mathbb Q[x]$, $g(x)\mapsto g(x+1)$,
jest automorfizmem pierścieni o odwrotności
$h(x)\mapsto h(x-1)$. Automorfizm zachowuje jednostki
i właściwe rozkłady; dlatego także $x^4+1$ jest
nierozkładalny nad $\mathbb Q$.
\end{przyklad}

\begin{twierdzenie}[Kryterium nierozkładalności przez redukcję modulo $p$]
\label{aa:tw:modp-irred}
Niech $f\in\mathbb Z[x]$ będzie pierwotnym wielomianem
dodatniego stopnia, a $p$ liczbą pierwszą, która nie
dzieli współczynnika wiodącego $f$. Jeśli wielomian
$\bar f\in\mathbb F_p[x]$ otrzymany przez redukcję
współczynników modulo $p$ jest nierozkładalny, to
$f$ jest nierozkładalny nad $\mathbb Q$.
\end{twierdzenie}
\begin{proof}
Gdyby $f$ był rozkładalny nad $\mathbb Q$, lemat
Gaussa dostarczyłby właściwego rozkładu
$f=gh$ z $g,h\in\mathbb Z[x]$ dodatnich stopni.
Współczynniki wiodące $g$ i $h$ mnożą się do
współczynnika wiodącego $f$, który nie jest podzielny
przez $p$. Żaden z tych dwóch współczynników nie jest
więc podzielny przez $p$. Po redukcji modulo $p$
stopnie $\bar g,\bar h$ pozostają dodatnie i nie zmieniają
się, a $\bar f=\bar g\bar h$ stanowi właściwy rozkład
w $\mathbb F_p[x]$. Sprzeczność dowodzi tezy.
\end{proof}

\begin{przyklad}[Kryterium redukcyjne jest jednostronne]
\label{aa:prz:modp-not-necessary}
Wielomian $x^2+1$ jest nierozkładalny nad $\mathbb Q$:
gdyby rozkładał się na czynniki liniowe, miałby pierwiastek
w $\mathbb Q$, a równanie $x^2=-1$ nie ma rozwiązania
nawet w $\mathbb R$. Tymczasem modulo $2$ mamy
$x^2+1=(x+1)^2$ w $\mathbb F_2[x]$.
Rozkładalność redukcji nie oznacza zatem
rozkładalności oryginalnego wielomianu.
\end{przyklad}

\begin{przyklad}[Nierozkładalność stopnia trzeciego wykazana modulo $2$]
\label{aa:prz:cubic-mod2-irreducibility}
Weźmy $f(x)=x^3+x+1\in\mathbb Z[x]$.
Współczynnik wiodący $1$ nie jest podzielny przez $2$.
Redukcja $\bar f\in\mathbb F_2[x]$ nie ma pierwiastków:
$\bar f(0)=1$, $\bar f(1)=1+1+1=1$.
Wielomian stopnia $3$ nad ciałem jest rozkładalny
wtedy i tylko wtedy, gdy ma czynnik stopnia $1$:
przy rozkładzie na dwa czynniki dodatniego stopnia
ich stopnie sumują się do $3$, więc jeden wynosi $1$.
Brak pierwiastków oznacza zatem, że $\bar f$
jest nierozkładalny. Kryterium redukcyjne daje
nierozkładalność $x^3+x+1$ nad $\mathbb Q$.
\end{przyklad}

\begin{przyklad}[Dlaczego trzeba zachować stopień przy redukcji]
\label{aa:prz:modp-leading-coefficient-necessary}
Wielomian $f(x)=2x^2+3x+1$ jest pierwotny:
jego współczynniki mają największy wspólny dzielnik
równy $1$. Mimo to $f$ rozkłada się nad
$\mathbb Q$ jako
$f(x)=(2x+1)(x+1)$, gdzie oba czynniki mają
stopień dodatni. Redukcja modulo $2$ daje natomiast
$\bar f(x)=x+1\in\mathbb F_2[x]$, który jest
nierozkładalny jako wielomian stopnia pierwszego.
Nie ma sprzeczności z
Twierdzeniem~\ref{aa:tw:modp-irred}: liczba $2$
dzieli współczynnik wiodący $f$, więc stopień
spada z dwóch do jednego. Po redukcji pierwszy
czynnik rozkładu staje się stałą $1$.
\end{przyklad}

\section{Pierwiastki wymierne i analiza konkretnych wielomianów}
\begin{twierdzenie}[Twierdzenie o pierwiastkach wymiernych]
\label{aa:tw:rational-root}
Niech $f=a_nx^n+\cdots+a_1x+a_0\in\mathbb Z[x]$
będzie wielomianem stopnia $n\ge1$ i niech $a_n\ne0$.
Jeśli $u/v\in\mathbb Q$ jest pierwiastkiem $f$,
gdzie $u\in\mathbb Z$, $v\in\mathbb N$ jest dodatnie
oraz $\gcd(u,v)=1$, to $u\mid a_0$ i $v\mid a_n$.
W przypadku $u=0$ pierwszy warunek oznacza $a_0=0$.
\end{twierdzenie}
\begin{proof}
Równość $f(u/v)=0$ po pomnożeniu przez $v^n$
daje
\[
 a_nu^n+a_{n-1}u^{n-1}v+\cdots+a_1uv^{n-1}+a_0v^n=0.
\]
Wyłączając pierwszy składnik, stwierdzamy, że
$v$ dzieli $a_nu^n$. Ponieważ $\gcd(u,v)=1$,
żaden element pierwszy dzielący $v$ nie dzieli $u$.
Lemat Euklidesa zastosowany wielokrotnie do potęg
czynników pierwszych $v$ daje $v\mid a_n$.
Analogicznie po wyłączeniu ostatniego składnika
otrzymujemy $u\mid a_0v^n$. Jeżeli $u\ne0$,
to $\gcd(u,v)=1$ daje $u\mid a_0$. Jeżeli $u=0$,
równość $f(0)=0$ oznacza $a_0=0$.
\end{proof}

\begin{przyklad}[Od listy kandydatów do pełnego rozkładu]
\label{aa:prz:rational-root-cubic}
Rozważmy $f(x)=2x^3-3x+1$. Każdy pierwiastek wymierny
w postaci nieskracalnej ma licznik dzielący $1$
i dodatni mianownik dzielący $2$. Kandydatami są
$\pm1$, $\pm\tfrac12$. Podstawienie $x=1$ daje
$f(1)=2-3+1=0$, więc z twierdzenia o czynniku liniowym
wiemy, że $x-1$ dzieli $f$.
Dzielenie wielomianów daje
$f(x)=(x-1)(2x^2+2x-1)$.
Równanie kwadratowe $2x^2+2x-1=0$ ma pierwiastki
$(-1\pm\sqrt3)/2$. Zatem otrzymaliśmy wszystkie
pierwiastki nad $\mathbb R$. Czynnika kwadratowego
nie można rozłożyć nad $\mathbb Q$, ponieważ
$\sqrt3\notin\mathbb Q$: gdyby $\sqrt3=m/n$
dla względnie pierwszych $m,n\in\mathbb Z$, $n\ne0$,
to $m^2=3n^2$. Stąd $3\mid m$, zatem po zapisaniu
$m=3k$ także $3\mid n$; byłaby to sprzeczność
z $\gcd(m,n)=1$.
\end{przyklad}

\begin{przyklad}[Brak pierwiastków wymiernych nie wystarcza w stopniu czwartym]
\label{aa:prz:rational-root-insufficient}
Wielomian $f(x)=x^4+4$ nie ma pierwiastków rzeczywistych,
ponieważ $x^4+4>0$ dla każdego $x\in\mathbb R$,
a więc nie ma też pierwiastków wymiernych.
Mimo to w $\mathbb Q[x]$ jest rozkładalny:
\[
 x^4+4=(x^2+2)^2-(2x)^2
       =(x^2-2x+2)(x^2+2x+2).
\]
Oba czynniki są niestałe. Twierdzenie o pierwiastkach
wymiernych pozwala więc wykluczyć rozkład stopnia $2$
lub $3$ (wtedy rozkład musi mieć czynnik liniowy),
lecz samo w sobie nie dowodzi nierozkładalności
wielomianów stopnia co najmniej $4$.
\end{przyklad}

\section{Pierścienie wielomianów wielu zmiennych i ich ilorazy}
\begin{definicja}[Pierścień wielomianów wielu zmiennych]
\label{aa:def:multi-polynomial}
Niech $R$ będzie przemiennym pierścieniem z jedynką
i niech $n\ge1$ będzie liczbą całkowitą.
Przyjmujemy $R[x_1]$ z Definicji~\ref{aa:def:poly-ring},
a dla $n\ge2$ definiujemy indukcyjnie
$R[x_1,\ldots,x_n]=(R[x_1,\ldots,x_{n-1}])[x_n]$.
Jego elementy są skończonymi sumami
$a_{\alpha}x_1^{\alpha_1}\cdots x_n^{\alpha_n}$,
gdzie $a_{\alpha}\in R$, a $\alpha_i\in\mathbb N_0$.
\end{definicja}

\begin{twierdzenie}[Pierścień współrzędnych paraboli]
\label{aa:tw:parabola-quotient}
Niech $K$ będzie ciałem. Homomorfizm pierścieni
$\psi:K[x,y]\to K[x]$, zadany wzorem
$\psi(F(x,y))=F(x,x^2)$, jest surjektywny, a jego
jądrem jest ideał $(y-x^2)$. W szczególności
$K[x,y]/(y-x^2)\cong K[x]$.
\end{twierdzenie}
\begin{proof}
Homomorfizm $\psi$ zachowuje sumy i iloczyny, ponieważ
jest podstawieniem, a jedynkę wysyła na jedynkę.
Jest surjektywny: każdy $h(x)\in K[x]$ jest obrazem
tego samego wielomianu traktowanego jako element
$K[x,y]$ niezależny od $y$.

Ponadto $\psi(y-x^2)=0$, więc ideał $(y-x^2)$ jest
zawarty w jądrze $\psi$. Aby udowodnić odwrotne
zawieranie, potraktujmy $F\in K[x,y]$ jako wielomian
zmiennej $y$ o współczynnikach z pierścienia $K[x]$.
Wielomian $y-x^2$ jest unormowany względem $y$.
Twierdzenie~\ref{aa:tw:poly-division} daje jednoznacznie
$F(x,y)=Q(x,y)(y-x^2)+H(x)$, ponieważ reszta ma
stopień względem $y$ mniejszy niż jeden.
Po podstawieniu $y=x^2$ otrzymujemy
$F(x,x^2)=H(x)$. Jeżeli $F\in\ker\psi$, to $H=0$
i $F\in(y-x^2)$. Mamy więc równość jądra i ideału.
Pierwsze twierdzenie o izomorfizmie pierścieni
(Rozdział~\ref{aa:ch:pierscienie}) daje wówczas
$K[x,y]/(y-x^2)\cong K[x]$.
\end{proof}

\begin{przyklad}[Równania opisujące ten sam zbiór nie muszą dawać tych samych pierścieni]
\label{aa:prz:poly-nilpotent-double}
W pierścieniu $\mathbb R[x]$ rozważmy ideały
$(x)$ oraz $(x^2)$. Zbiory rzeczywistych zer wielomianów
$x$ i $x^2$ są identyczne: w obu przypadkach składają
się z jednego punktu $0$. Jednak
$\mathbb R[x]/(x)\cong\mathbb R$ jest ciałem,
natomiast w $\mathbb R[x]/(x^2)$ klasa $x+(x^2)$ jest
niezerowa i nilpotentna stopnia dwa.
Niezerowość tej klasy wynika z faktu, że $x$ nie jest
wielokrotnością $x^2$. Oba pierścienie nie mogą być
izomorficzne: izomorfizm zachowywałby niezerowe elementy
nilpotentne, których w ciele nie ma. W następnym rozdziale
wyjaśnimy ten przykład za pomocą pierwiastnika ideału.
\end{przyklad}

\begin{uwaga}[Przejście do teorii Galois]
\label{aa:uw:poly-galois}
Kryteria nierozkładalności rozstrzygają, kiedy iloraz
$K[x]/(f)$ jest ciałem: na mocy
Twierdzenia~\ref{aa:tw:prime-max-quotient} zachodzi to
dokładnie wtedy, gdy ideał $(f)$ jest maksymalny.
W $K[x]$ nierozkładalność wielomianu $f$ dodatniego
stopnia jest równoważna maksymalności ideału $(f)$;
wynika to z algorytmu Euklidesa, który daje
odwrotność klasy każdego wielomianu $g$ niepodzielnego
przez $f$. Ta obserwacja pozwoli konstruować
rozszerzenia ciał i przejść do teorii Galois.
\end{uwaga}
